\section{The complete argumentation meta-interpreter construction}
\label{sec:annex-p18-argumentation-meta-interpreters}

\begin{annexitem}{annex:p18.meta.reduct}{thm:p18-meta-reduct}{Every reduct atom and the defense interface}
\begin{proof}
In \(Q^X\), the acceptance clause for \(a\) is deleted exactly
when \(a^\bullet\in X\); otherwise it becomes the fact
\(a^\circ\leftarrow\). The defeat clauses have no negative
premises and remain unchanged. Starting from the empty set,
positive consequence therefore first derives exactly the
acceptance atoms for \(A\setminus d_X\), and then exactly
the defeat atoms for \((A\setminus d_X)^+\).
No defeat atom can generate a further acceptance atom in
this positive reduct, because its acceptance rules are
already facts or deleted. The resulting set is closed
under all clauses, and every model contains the facts and
their defeat consequences. It is the least model, proving
the formula for \(\Gamma_Q\).

For \(P_D^X\), the extensional facts are never deleted and
are the only rules for their predicates. They first give
exactly \(B_D^+\). A remaining acceptance rule is now
\(\operatorname{acc}(a)\leftarrow\operatorname{arg}(a)\);
it derives its head exactly for \(a\in A\setminus d_X\).
For names outside \(A\), the domain premise is absent from
the least model, so they give no acceptance head. A defeat
rule derives its head exactly when both names are in \(A\),
the stated edge is a fact, and its source acceptance atom
has been derived. These are exactly the targets
\((A\setminus d_X)^+\). This set together with the preceding
facts is closed, and every positive model contains it.
There is no other predicate or defining clause, so the
formula accounts for every atom, including extra names
and nonedges. It remains valid if \(A\) is empty.

For the literal source program \(P^0_{\mathcal F}\), all
ground names are actual arguments and there are no domain
predicates. Its reduct keeps the attack facts, makes
acceptance a fact exactly on \(A\setminus d_X\), and
derives defeat exactly at their attack targets.
This is \(m_0(A\setminus d_X)\), proving the source
formula by its own least-model calculation.

For \(X=m(S)\), its defeat component is \(d_X=S^+\).
Substitution gives \(\Gamma(m(S))=m(N(S))\).
The no-attack operator satisfies
\[
\begin{aligned}
 a\in N(N(S))
 &\Longleftrightarrow
       \text{no attacker of }a\text{ belongs to }A\setminus S^+\\
 &\Longleftrightarrow \operatorname{Att}(a)\subseteq S^+
 \Longleftrightarrow a\in\mathsf F(S).
\end{aligned}
\]
Applying the first identity twice proves the second.
The least-model calculations use finite bodies and actual
edge witnesses, not a finite bound on the number of
clauses or attackers.
\end{proof}
\end{annexitem}

\begin{annexitem}{annex:p18.meta.stable}{thm:p18-meta-stable}{The two inverse stable-model maps}
\begin{proof}
Suppose \(E\) is stable in \(\mathcal F\).
By \cref{thm:p18-arbitrary-admissibility}, stability is
exactly \(E=A\setminus E^+\). The reduct identity gives
\(\Gamma(m(E))=m(E)\), so \(m(E)\) is a stable model
for each of the three stated encodings at its declared domain.

Conversely let \(M=\Gamma(M)\) be a stable model.
The complete reduct calculation forces \(M=m(E)\),
where \(E=A\setminus d_M\). Reading its defeat component
gives \(d_M=E^+\); hence \(E=A\setminus E^+\).
This states both conflict-freeness and attack of every
outside argument, so \(E\) is stable.
The acceptance component of \(m(E)\) is precisely \(E\).
It recovers \(E\) uniquely, proves injectivity, and proves
that the displayed constructions and extractions are
inverse in both directions.

For stable \(E\), the sets \(E,E^+\) are disjoint and
cover \(A\). Thus every dynamic atom receives exactly
one sign in \(L_Q(E)\): acceptance is true on \(E\)
and false on \(E^+\), while defeat has the opposite
truth value. The lift additionally assigns the correct
sign to every extensional atom and every atom at an
extra name. Hence \(L_D(E)\) is the total interpretation
of \(m_D(E)\). For the empty framework, \(Q\) has no
atoms and its sole stable model is empty; for \(P_D\),
the same calculation gives the empty positive model
and the total negative background specified in \(B_D\).
Deleting the domain literals when \(A\ne\varnothing\)
gives the asserted total interpretation \(L_0(E)\)
for the source program.
\end{proof}
\end{annexitem}

\begin{annexitem}{annex:p18.meta.specialization}{thm:p18-meta-specialization}{The complete extensional and grounding bridge}
\begin{proof}
Consider a fixed point \(J=\Phi_{P_D}(J)\).
An extensional fact has a true empty body and no other
defining clauses, so it is true and cannot be false.
An extensional atom which is not a fact has no defining
clauses, so all its bodies are false vacuously and it
is false, never true. Thus all the extensional signs
in \(B_D\) are forced, and their opposite signs are
absent.

If \(d\notin A\), the atom \(\operatorname{arg}(d)\)
is false. The acceptance clause for \(d\) has a false
body, and every defeat clause for \(d\) also has that
false domain premise. Hence both dynamic atoms at
\(d\) are false, never true, at a fixed point.
This proves every part of the fixed background \(B_D\).

Now fix \(a\in A\). Its domain atom is true.
The acceptance clause has the same truth value as
\(\neg a^\bullet\). For a ground defeat clause with
target \(a\), an instance with source outside \(A\)
has a false source-domain premise; an instance with
source \(b\in A\) but \((b,a)\notin R\) has a false
attack premise. For an actual edge, all three
extensional premises are true, so the body's value
is exactly the value of \(b^\circ\).
It follows that the two dynamic fixed-point equations
are exactly those of
\[
          a^\circ\leftarrow\neg a^\bullet,\qquad
          a^\bullet\leftarrow b^\circ\quad(b\to a).
\]
In particular no incoming edge is a universal
condition over an omitted finite list: falsity of
the defeat head requires falsity of every actual
incoming acceptance premise.

Consequently restricting \(J\) to \(H_{\mathcal F}\)
gives a fixed point \(I=\Phi_Q(I)\), and
\(J=\iota_D(I)\). Conversely the same truth, falsity
and empty-definition calculations show that any
\(\iota_D(I)\) with \(I=\Phi_Q(I)\) is fixed for \(P_D\).
Consistency is preserved and reflected, because the
background is consistent and uses no dynamic atom
of \(A\). Both extension by a fixed background and
restriction preserve and reflect inclusion.

The least Fitting fixed points exist and are consistent
by the previously proved \cref{thm:p18-clause-fitting},
which allows arbitrary ground programs with finite
bodies. Their order-isomorphic fixed-point families
therefore have corresponding least elements. This
proves the displayed Fitting equality. Equivalently,
the least fixed point of \(P_D\) lies below the lift
of the least fixed point of \(Q\); restricting it
gives the opposite leastness inequality.
For the source program on \(D=A\), every domain
premise above is already true by its specified
quantifier domain; deleting that predicate and its
literals gives exactly the same proof.
\end{proof}
\end{annexitem}

\begin{annexitem}{annex:p18.meta.defense.steps}{thm:p18-meta-defense-steps}{All four successor steps and every limit stage}
\begin{proof}
Represent a dynamic literal set \(I\) by four subsets
of \(A\):
\[
\begin{aligned}
 T&=\{a:a^\circ\in I\},&
 Z&=\{a:\neg a^\circ\in I\},\\
 D&=\{a:a^\bullet\in I\},&
 U&=\{a:\neg a^\bullet\in I\}.
\end{aligned}
\]
These names describe truth and falsity components;
they impose no consistency condition in this auxiliary
calculation. Put
\[
                  B(Z)=\{a:\operatorname{Att}(a)\subseteq Z\}.
\]
The acceptance rule is true when its defeat atom is
false and false when that atom is true. The defeat
head is true if some attacker is accepted, and false
if every attacker is rejected. Thus
\[
                          \Phi_Q(T,Z,D,U)=(U,D,T^+,B(Z)).
\]
This includes an unattacked head: its defeat definition
is empty, so it is false.

The four components of \(L_Q(S)\) are
\((S,S^+,S^+,S)\), and \(B(S^+)=\mathsf F(S)\).
Writing \(F=\mathsf F(S)\) only in the following
calculation, the successive quadruples are
\[
\begin{array}{c|cccc}
 &\text{accept true}&\text{accept false}
       &\text{defeat true}&\text{defeat false}\\ \hline
0&S&S^+&S^+&S\\
1&S&S^+&S^+&F\\
2&F&S^+&S^+&F\\
3&F&S^+&F^+&F\\
4&F&F^+&F^+&F .
\end{array}
\]
The last row is \(L_Q(\mathsf F(S))\), proving the
exact four-step identity. If \(S=\mathsf F(S)\),
already the first row is fixed by \(\Phi_Q\).

For the transfinite assertion, stage zero is
\(L_Q(\varnothing)=\varnothing\). At a successor,
the established four-step identity advances the
program stage from \(\tau_\alpha\) to
\(\tau_\alpha+4=\tau_{\alpha+1}\).
At a nonzero limit \(\lambda\), the stages
\(\tau_\beta\), \(\beta<\lambda\), are cofinal
in \(\tau_\lambda\).
The Fitting iteration is increasing from its empty
initial set, by monotonicity, so its union at
\(\tau_\lambda\) equals the union over this cofinal
subsequence. Furthermore attack images commute with
unions:
\[
              \left(\bigcup_{\beta<\lambda}S_\beta\right)^+
                      =\bigcup_{\beta<\lambda}S_\beta^+.
\]
Indeed an attack has one actual source belonging to
some member of the union. Every component of \(L_Q\)
therefore commutes with this union. These facts give
the required limit equality and complete the induction.

The argument iteration reaches the grounded extension
at a set-sized ordinal by
\cref{thm:p18-transfinite-grounded}.
At that program stage the value is \(L_Q(G)\),
and the first-row calculation shows it is fixed.
An increasing iteration from the empty set is below
every fixed point at every stage, by transfinite
induction. Hence this reached fixed point is the
least one, proving \(\operatorname{FM}(Q)=L_Q(G)\).
\end{proof}
\end{annexitem}

\begin{annexitem}{annex:p18.meta.completion}{thm:p18-meta-completion}{Every partial model and all extension semantics}
\begin{proof}
Use the four components \((T,Z,D,U)\) from the preceding
proof. At a fixed point they satisfy
\[
                       T=U,\quad Z=D,\quad D=T^+,\quad U=B(Z).
\]
Consequently \(Z=D=T^+\) and
\(T=B(T^+)=\mathsf F(T)\).
Consistency requires \(T\cap Z=\varnothing\) and
\(D\cap U=\varnothing\); with these equations both
conditions are exactly \(T\cap T^+=\varnothing\).
Thus \(T\) is a conflict-free fixed point of defense,
which is exactly a complete extension, and the full
literal set is \(L_Q(T)\).

Conversely, for complete \(E\), conflict-freeness makes
\(L_Q(E)\) consistent and \(\mathsf F(E)=E\) makes all
four displayed equations hold. This proves the two
inverse maps. If \(E\subseteq E'\), then
\(E^+\subseteq E'^+\), so every component of \(L_Q(E)\)
is contained in its counterpart in \(L_Q(E')\).
Conversely inclusion of literal sets implies inclusion
of their true acceptance components. The maps therefore
preserve and reflect order.
The lift theorem transfers this isomorphism to \(P_D\).

The grounded extension is the least complete one,
so it corresponds to the least partial completion
model. Preferred extensions are exactly the maximal
complete extensions: a preferred set is complete,
and a maximal complete set cannot have a proper
admissible superset, since such a superset would
extend to a larger preferred, hence complete, set
by \cref{thm:p18-arbitrary-admissibility}.
The isomorphism transfers maximality in both
directions.

Finally \(L_Q(E)\) is total precisely when
\(E\cup E^+=A\), since these are exactly the arguments
receiving signs for their dynamic atoms. Its
consistency already gives disjointness. Totality
is therefore exactly stability of \(E\).
The background is total on every other atom, so
the same assertion holds after lifting. This proves
the total-completion statement and also explains
why the total completion models of this particular
encoding are its stable models. The derivation uses
the pointwise fixed-point definition for all ground
heads; no finite bound on an incoming disjunction
has been assumed.
\end{proof}
\end{annexitem}

\begin{annexitem}{annex:p18.meta.positive.rank}{thm:p18-meta-positive-rank}{The full unfounded-set calculation by positive rank}
\begin{proof}
Use \(Z_P(I)\) for the atoms whose every defining body
is false, and \(U_P(I)\) for the greatest unfounded set.
Always \(Z_P(I)\subseteq U_P(I)\): the set \(Z_P(I)\)
is itself unfounded, since all relevant bodies are
false. This includes atoms with no defining clauses.

Suppose \(I=\Phi_P(I)\). Then \(I^-=Z_P(I)\).
We prove that every member of \(U_P(I)\) belongs
to \(I^-\), by induction on its finite rank.
For a rank-zero atom in \(U_P(I)\), no clause for
it has a positive body atom, because such a premise
would require a smaller rank. Unfoundedness therefore
requires every defining body to be false; the atom
lies in \(Z_P(I)=I^-\).

At a larger rank, inspect any clause whose head
is the selected unfounded atom. Its body is either
false already or has a positive premise in \(U_P(I)\).
That premise has strictly smaller rank and therefore,
by induction, belongs to \(I^-\). Its negative
literal in \(I\) makes the clause's body false.
Thus every defining body is false in either case,
putting the head in \(Z_P(I)=I^-\).
This argument quantifies over every defining clause,
including infinitely many of them. We obtain
\[
                          U_P(I)=Z_P(I)=I^-.
\]
Since \(\Phi_P(I)=I\) also gives \(T_P(I)=I^+\),
it follows that \(V_P(I)=I\).

Conversely suppose \(I=V_P(I)\). Then
\(I^-=U_P(I)\). For any head in this unfounded set,
each defining clause is false or has a positive
premise in \(U_P(I)=I^-\). In the latter case
that premise is false in \(I\), so again the body
is false. Thus \(U_P(I)\subseteq Z_P(I)\).
The reverse inclusion always holds, giving equality,
and \(T_P(I)=I^+\) now implies \(\Phi_P(I)=I\).
This converse does not require the rank hypothesis.

The argument is valid on the auxiliary full literal
powerset, even before consistency is imposed. Both
operators are monotone and have least fixed points,
as proved in the program-semantics foundation.
Equality of their fixed-point sets proves equality
of their least fixed points.
Each stated interpreter ranking strictly decreases
on every positive premise: the reduced defeat
rules depend positively only on acceptance; the
guarded acceptance rules depend positively on
domain facts, and guarded defeat rules on domain,
attack and acceptance. Extensional facts have empty
bodies. Negative defeat premises impose no positive
rank condition. This verifies all three applications
without requiring the original attack relation
to be acyclic.
\end{proof}
\end{annexitem}

\begin{annexitem}{annex:p18.meta.full.correspondence}{thm:p18-meta-full-correspondence}{The source theorem with every output literal and its graph bridges}
\begin{proof}
The four-step construction gives
\(\operatorname{FM}(Q)=L_Q(G)\).
The specialization theorem gives
\(\operatorname{FM}(P_D)=\iota_D(L_Q(G))=L_D(G)\).
The rank theorem identifies each of these Fitting
models with the corresponding well-founded model.
This proves the forward grounded implications.
For the source program the specialization theorem with
domain literals omitted gives
\(\operatorname{FM}(P^0_{\mathcal F})=L_0(G)\).
Its positive rank then identifies this with its
well-founded model as well.

Conversely if \(L_Q(E)\), \(L_D(E)\), or, in the
source domain, \(L_0(E)\), equals the
appropriate well-founded model, extract its true
acceptance atoms. The true acceptance set of the
displayed model is \(G\), and that of the signed
encoding is exactly \(E\). Therefore \(E=G\).
This proves the required reverse implications for
every proposed \(E\subseteq A\).

Grounded conflict-freeness gives \(G\cap G^+=\varnothing\).
For a member of \(G\), the definition of \(L\) gives
true acceptance and false defeat, and for a member
of \(G^+\) it gives true defeat and false acceptance.
For a member outside both sets it gives neither
sign for either atom, which is precisely undefinedness.
The complete fixed background was proved by the
grounding calculation. Thus no false dynamic,
domain or nonedge literal is omitted.
The earlier stable-model theorem supplies both
directions for the positive map \(m\); together
these statements establish the three claimed
interpreter correspondences.

For the source's program-to-argumentation proof
route, apply the already proved
\cref{thm:p18-program-well-founded} to the actual
ground program \(P_D\). Every one of its rule
bodies is finite, and its specified atom base
is a set, so the arbitrary-ground-program
premises of that theorem apply. Its full
Appendix~B interface, \cref{thm:p18-unfounded-support},
therefore gives
\[
 \operatorname{con}(G_{\mathcal A_\infty(P_D)})
                         =\operatorname{WFM}(P_D)=L_D(G).
\]
Here the graph is the literal support graph:
all finite supports, including unused assumptions,
and all negative unit arguments are retained.
No restriction of the attack clause to defeat
conclusions has been made.
Similarly \cref{thm:p18-clause-fitting} applies
to the literal clause graph and gives
\[
 \operatorname{con}(G_{\mathcal A_{\rm c}(P_D)})
                         =\operatorname{FM}(P_D)=L_D(G).
\]
These equalities identify conclusions, not the
three different argument universes.

There is also a complete stable bridge through
the source support graph. For stable \(E\),
\(M=m_D(E)\) is a stable program model. With
\(H_{P_D}\) the full program base,
\cref{thm:p18-program-stable-correspondence}
constructs its unique stable support extension
\[
                    {\cal E}_\infty(H_{P_D}\setminus M).
\]
Conversely any stable extension of that graph
gives a stable program model by its positive
conclusions, and the true acceptance atoms of
that model give a unique stable \(E\) in
\(\mathcal F\). This states both directions
without replacing the source program graph
by a similar smaller argument system.
\end{proof}
\end{annexitem}

\begin{annexitem}{annex:p18.meta.finite.procedure}{thm:p18-meta-finite-procedure}{Finite construction, closure and stable enumeration}
\begin{proof}
There is one reduced acceptance clause for each
listed argument and one reduced defeat clause
for each listed edge. Each has constant size
apart from its names, giving \(n+m\) clauses
and the declared \(2n\)-atom base. Isolated
arguments still receive acceptance clauses;
the empty input produces the empty program.

For a defense evaluation, the first edge scan
collects exactly \(S^+\), by the definition of
the attack image. In the second scan, an
argument is rejected exactly if it has an
attacker outside \(S^+\). The surviving set
is therefore exactly \(\mathsf F(S)\).
Initialization and the two scans cost
\(O(n+m+1)\), including the empty input. The finite grounded theorem,
\cref{thm:p18-grounded-iteration}, proves
termination with the grounded answer after
at most \(n+1\) evaluations, including the
final equality check. Its stated time bound
follows. The signed output is correct by
the full correspondence theorem.

For stable enumeration, each candidate \(E\)
has its image \(E^+\) computed by an edge
scan. The two tested conditions are exactly
internal conflict-freeness and attack of
every outside argument. Thus the retained
sets are all and only stable extensions.
The reduct bijection gives exactly all
stable program models, with neither omissions
nor duplicate candidates. There are \(2^n\)
subsets and each check costs \(O(n+m+1)\),
including the empty input. This proves the
enumeration bound.

If the finite domain is \(D=A\), there
are \(n^2\) possible attack atoms; listing
the signs of every nonedge can therefore
require quadratic background output even
when the dynamic program is sparse.
For a larger or infinite \(D\), the full
background may itself be infinite.
No such output is covered by a finite
dynamic-output space bound.

The source's two-rule schema plus extensional
facts can be finitely supplied when \(A,R\)
are finite. General infinite domains require
a representation of the argument and attack
data. An enumeration of positive facts
does not certify that all attackers of a
given point have been found. Thus the
finite closure algorithm's universal
defense step and terminating stable
enumeration do not follow from the
semantic theorem for infinite inputs.
\end{proof}
\end{annexitem}

\begin{annexitem}{annex:p18.meta.boundary.examples}{thm:p18-meta-boundary-examples}{All stated boundary and semantics counterexamples}
\begin{proof}
For the one-edge framework \(a\to b\), the argument
\(a\) is unattacked, so it belongs to \(\mathsf F(\varnothing)\).
The argument \(b\) is attacked by \(a\), and no argument
attacks \(a\); hence \(b\) is not defended by any set.
The grounded extension is therefore \(\{a\}\).
Any stable extension must contain the unattacked
\(a\), and then conflict-freeness excludes \(b\);
the resulting set attacks its only outsider, so
it is the unique stable extension.
Its image is \(\{b\}\), giving exactly the four
dynamic literals displayed in the statement.
In particular the rule
\(b^\circ\leftarrow\neg b^\bullet\) has a false
body when \(b^\bullet\) is true. Every Fitting
or well-founded fixed point therefore includes
\(\neg b^\circ\); omitting it cannot specify
the full model.

For an isolated \(a\), its defeat atom has no
defining reduced clause and is false. Its
acceptance rule therefore gives true acceptance.
This yields the displayed two-literal answer.
There are no attack facts mentioning the name,
so a construction which takes only names found
in those facts loses the acceptance clause
and cannot represent this answer.

For an extra name \(d\notin A\), the guarded
defeat rules have a false target-domain premise,
so \(\operatorname{defeat}(d)\) is false.
With the acceptance guard it also has a false
acceptance body, so acceptance is false.
If that guard is removed, its acceptance body
is now the true negative literal
\(\neg\operatorname{defeat}(d)\), and acceptance
becomes true. The same calculation with
\(A=\varnothing,D=\{d\}\) proves the empty-framework
counterexample. In the entirely unguarded source
schema over that one-name domain, defeat has no
true edge body and the same false-defeat,
true-acceptance result follows.

For the self-loop \(a\to a\), the reduced program is
\[
                    a^\circ\leftarrow\neg a^\bullet,\qquad
                    a^\bullet\leftarrow a^\circ.
\]
The empty literal set is a Fitting fixed point:
the negative acceptance body and the positive
defeat body are both undefined. Thus both
Fitting and well-founded models are dynamically
empty. The only framework subsets are empty,
which fails to attack its outsider, and
\(\{a\}\), which is not conflict-free.
There is no stable extension and therefore
no stable program model. For the mutual attack
\(a\leftrightarrow b\), no argument is defended
by the empty set, so the grounded set is empty.
Each singleton is stable. The empty set is
not stable and the pair is conflicting, so
these are the only two stable sets.
Their positive program models are respectively
\(\{a^\circ,b^\bullet\}\) and
\(\{b^\circ,a^\bullet\}\); the dynamic least
partial model assigns no sign to any of the
four atoms.

For an isolated \(a\), ordinary material
implication permits defeat to be true without
support: assigning \(a^\bullet\) true makes
the antecedent of
\(\neg a^\bullet\mathbin{\to}a^\circ\) false.
There is no incoming defeat rule to violate.
With false acceptance this is a classical
model of the implications. It is not the
least model of its reduct, and its true
defeat atom violates the empty completed
defeat definition. Hence it is neither a
stable nor a completion model.

For the predicate mismatch, stipulate that
\(\operatorname{attacks}\) and
\(\operatorname{atk}\) are distinct symbols.
The sole supplied edge fact then defines
no \(\operatorname{atk}\) atom. Every
processing-rule edge body is false, both
defeat atoms are false, and both arguments
are accepted. The accepted pair conflicts
in the intended one-edge framework.
This proves the signature counterexample.

Finally the guarded program for any actual
argument \(a\) has a finite support proof of
\(a^\circ\) using its domain fact and the
negative assumption \(\neg a^\bullet\).
The corresponding argument
\((\{\neg a^\bullet\},a^\circ)\) is present
in the literal support graph. That graph
also contains the unit
\((\{\neg a^\circ\},\neg a^\circ)\).
The complement of the first conclusion is
in the second support, so the former attacks
the latter by the exact source attack clause.
Its conclusion is acceptance, not defeat.
This also shows directly why retaining
all unit arguments matters for the full
signed semantics.
\end{proof}
\end{annexitem}

\begin{annexitem}{annex:p18.meta.infinite.input}{thm:p18-meta-infinite-input}{A decidable graph whose unique answer encodes nonhalting}
\begin{proof}
Fix an effective enumeration of Turing machines
run on the empty input. Use the decidable
argument-name set
\[
           A=\{a_e:e\in\mathbb N\}
                       \cup\{b_{e,t}:e,t\in\mathbb N\}.
\]
There is an attack \(b_{e,t}\to a_e\) exactly
when machine \(e\) halts within \(t\) steps.
There are no other attacks. The edge predicate
is decidable by a bounded simulation after
checking the two name forms. Every path has
length at most one, since no \(a_e\) attacks
anything and no \(b_{e,t}\) is attacked.

Every \(b_{e,t}\) is unattacked and belongs
to the first defense stage.
If machine \(e\) never halts, \(a_e\) also
has no attacker and belongs to that stage.
If it halts, \(a_e\) has an attacker
\(b_{e,t}\) which no argument can attack.
It is therefore defended by no set and
never belongs to a later stage.
Consequently
\[
 G=\{b_{e,t}:e,t\in\mathbb N\}
                      \cup\{a_e:e\text{ does not halt}\}.
\]
This set is conflict-free, and every outsider
is some \(a_e\) for a halting machine and is
attacked by an included \(b_{e,t}\).
It is stable. Any stable set must contain
every unattacked argument, and including
all \(b_{e,t}\) excludes every halting
\(a_e\) by conflict-freeness. Thus \(G\)
is the unique stable extension.

If grounded membership were decidable for
these effectively named inputs, testing
membership of \(a_e\) would decide nonhalting
and hence decide halting, contradicting the
halting theorem. The undecidability premise here
has the usual direct diagonal proof: if a
procedure decided whether program \(e\) halts
on input \(x\), a program on input \(x\) could
loop when that procedure says that \(x\)
halts on itself, and halt otherwise.
Running this program on its own code
contradicts either answer. A decider for
empty-input halting would decide the
general problem by compiling a program
which ignores its input and simulates
the specified program on the specified
input. Thus the empty-input decision
used in this construction is impossible.
Since \(G\cup G^+=A\),
its signed dynamic model \(L_Q(G)\) is total.
The interpreter correspondences make it
both the well-founded and Fitting model,
and its positive part the unique stable
model. This proves every assertion while
keeping the attack relation itself decidable.
\end{proof}
\end{annexitem}
